\documentclass[12pt,a4paper]{article}
\usepackage{graphicx}
\usepackage{booktabs}
\usepackage{amsmath}
\usepackage{caption}
\usepackage{geometry}
\usepackage{array}
\geometry{margin=2.5cm}

\renewcommand{\thefigure}{S\arabic{figure}}
\renewcommand{\thetable}{S\arabic{table}}

\title{\textbf{Supplementary Material}\\[0.5em]
\large Frontier-Centered Local Map Completion for Autonomous Exploration}
\author{Selma G\"{u}zel \and Erkan Uslu}
\date{}

\begin{document}

\maketitle

This document provides supplementary results referenced in the main paper,
including qualitative model comparisons (Figs.~S1--S2), UNet variant
evaluation (Table~S1), and dataset size sensitivity results (Table~S2).

% ============================================================
\section*{Dataset Split Details}
% ============================================================

The dataset was initially partitioned into train (81.4\%),
validation (8.6\%), and test (10\%) sets at the environment
level, yielding 9,418 train, 999 validation, and 1,160 test
patches. To construct the larger dataset (BDS) without
additional data collection, 1,000 patches were moved from
the original test set into training, resulting in 10,418
train patches. The remaining 160 test patches form the
shared test set used for evaluation of both SDS and BDS
configurations. The validation set is identical across
both configurations.

% ============================================================
\section*{Qualitative Model Comparisons}
% ============================================================

Figures~\ref{fig:qualitative_1} and~\ref{fig:qualitative_2} show
qualitative predictions of all evaluated models on the SDS test set.
\texttt{UNet\textsubscript{base}} preserves thin walls and topological
continuity most consistently. GAN variants (Pix2Pix-hinge, Seg-GAN)
produce occasional hallucinations near obstacle boundaries. ViT
generates fragmented wall predictions due to coarse $16\times16$ patch
tokenisation. CVAE fails to capture structural patterns under irregular
frontier masks. LSTM and ImageGPT produce near-uniform or structurally
incoherent outputs.

\begin{figure}[!ht]
    \centering
    \includegraphics[width=\textwidth]{frontier_qualitative_all_models_1.png}
    \caption{Qualitative comparison of frontier-centered completion
    models on SDS. Columns left to right: input condition, ViT, UNet,
    Pix2Pix, GEOM, CVAE, blended (known cells from input, predicted
    cells from model), target.}
    \label{fig:qualitative_1}
\end{figure}

\begin{figure}[!ht]
    \centering
    \includegraphics[width=\textwidth]{frontier_qualitative_all_models_2.png}
    \caption{Qualitative comparison of frontier-centered completion
    models on SDS. Columns left to right: input condition, SegGAN,
    LSTM, ImageGPT, blended (known cells from input, predicted cells
    from model), target.}
    \label{fig:qualitative_2}
\end{figure}

\clearpage

% ============================================================
\section*{Table S1: UNet Variant Comparison (MASKED, SDS)}
% ============================================================

Table~\ref{tab:unet_variants} reports MASKED evaluation results for
all U-Net architectural variants on SDS. Since IoU and F1 differences
are small (IoU range: 0.316--0.334), model selection is based on the
average of all metrics (Pixel Acc, F1, IoU, L1, normalised throughput).
\texttt{UNet\textsubscript{base}} achieves the highest overall average
(0.627), driven by its superior Pixel Accuracy (0.649) and L1 (0.192),
despite slightly lower IoU than deeper variants. LaMa, while a
Fourier-augmented U-Net variant, is excluded from online deployment
due to its throughput (12.67~img/s).

\begin{table}[!ht]
\renewcommand{\arraystretch}{1.3}
\caption{UNet Variant Comparison: MASKED Evaluation on SDS.
Selection criterion: average of all metrics (Pixel Acc, F1, IoU,
L1, normalised throughput). $\dagger$ = excluded (throughput $<$ 40~img/s).}
\label{tab:unet_variants}
\centering
\setlength{\tabcolsep}{6pt}
\begin{tabular}{lccccc}
\toprule
\textbf{Model} & \textbf{Pixel Acc}$\uparrow$
    & \textbf{F1-mac}$\uparrow$
    & \textbf{IoU-mac}$\uparrow$
    & \textbf{L1}$\downarrow$
    & \textbf{img/s}$\uparrow$ \\
\midrule
\textbf{UNet\textsubscript{base}} & \textbf{0.649} & 0.412 & 0.316 & \textbf{0.192} & 42.25 \\
UNet\textsubscript{new}           & 0.589 & 0.436 & 0.329 & 0.231 & \textbf{44.47} \\
UNet\textsubscript{att}           & 0.587 & \textbf{0.441} & 0.333 & 0.228 & 39.33 \\
UNet (4 layers)                   & 0.562 & \textbf{0.441} & \textbf{0.334} & 0.237 & 43.63 \\
LaMa$^\dagger$                    & 0.541 & 0.428 & 0.318 & 0.250 & 12.67 \\
\bottomrule
\end{tabular}
\end{table}

\clearpage

% ============================================================
\section*{Table S2: Effect of Dataset Size (SDS $\rightarrow$ BDS)}
% ============================================================

Table~\ref{tab:masked_bds} reports MASKED results when models are
retrained on the larger dataset (BDS: 10,418 training patches vs.\
SDS: 9,418). \texttt{UNet\textsubscript{base}} performance remains
nearly unchanged (IoU 0.309 vs.\ 0.316 on SDS), confirming the
geometry-dominant nature of frontier-centered completion. GAN variants
benefit modestly; ViT improves but remains below U-Net. This
stability may not hold for substantially larger datasets, where
data-hungry architectures such as ViT could potentially close the gap.
RePaint and GEOM are omitted as they fall below the 40~img/s threshold
or produce near-zero MASKED scores on SDS.

\begin{table}[!ht]
\renewcommand{\arraystretch}{1.3}
\caption{MASKED Evaluation on BDS. U-Net performance is stable
across dataset sizes; other families benefit more from additional data.}
\label{tab:masked_bds}
\centering
\setlength{\tabcolsep}{6pt}
\begin{tabular}{lcccc}
\toprule
\textbf{Model} & \textbf{Pixel Acc}$\uparrow$
    & \textbf{F1-mac}$\uparrow$
    & \textbf{IoU-mac}$\uparrow$
    & \textbf{L1}$\downarrow$ \\
\midrule
\textbf{UNet\textsubscript{base}} & \textbf{0.619} & 0.409 & 0.309 & \textbf{0.206} \\
UNet\textsubscript{new}           & 0.566 & \textbf{0.437} & \textbf{0.329} & 0.238 \\
Pix2Pix-hinge                     & $\approx$0.60 & $\approx$0.39 & $\approx$0.30 & $\approx$0.21 \\
ViT                               & 0.570 & 0.323 & 0.242 & 0.236 \\
CVAE                              & 0.404 & 0.267 & 0.180 & 0.317 \\
\bottomrule
\end{tabular}
\end{table}

\end{document}
